% Generated draft: requires mathematical and source review.
\documentclass[11pt]{article}
\usepackage[margin=25mm]{geometry}
\usepackage{fontspec}
\setmainfont{Latin Modern Roman}
\usepackage{amsmath,amssymb,mathtools,unicode-math}
\setmathfont{Latin Modern Math}
\usepackage{enumitem,needspace,xurl,hyperref}
\setlength{\parindent}{0pt}
\setlength{\parskip}{5pt}
\setlength{\emergencystretch}{3em}
\newcommand{\N}{\mathbb{N}}
\newcommand{\Z}{\mathbb{Z}}
\newcommand{\Q}{\mathbb{Q}}
\newcommand{\R}{\mathbb{R}}
\newcommand{\C}{\mathbb{C}}
\newcommand{\F}{\mathbb{F}}
\newcommand{\A}{\mathbb{A}}
\newcommand{\PP}{\mathbb{P}}
\newcommand{\E}{\mathbb{E}}
\newcommand{\eps}{\varepsilon}
\newcommand{\dd}{\,\mathrm{d}}
\begin{document}
\section*{UnsolvedMath ID: 30005780}
\textbf{Source problem code:} catalogue entry
\subsection{Definitions and mathematical statement}
Let \(k\) be a field and let \(A\) be a finite-dimensional Hopf algebra over \(k\), without assuming cocommutativity. Give \(k\) its trivial \(A\)-module structure through the counit, and set
\[
R_A=H^*(A,k)=\bigoplus_{n\geq 0}\operatorname{Ext}^n_A(k,k),
\]
with the Yoneda product and cohomological grading. No characteristic or algebraic-closure assumption on \(k\) is supplied; the questions are therefore parametrized by \((k,A)\), and any answer must state its field hypotheses.

Let \(\operatorname{stab}A\) denote the stable category of finite-dimensional left \(A\)-modules: morphisms are module homomorphisms modulo those factoring through projective modules. It is a triangulated category with tensor product over \(k\), where \(A\) acts through its coproduct. In a symmetric tensor triangulated setting, the Balmer spectrum consists of proper prime thick tensor ideals, with its usual topology. Here a thick subcategory is triangulated and closed under direct summands; a tensor ideal is closed under tensoring with arbitrary objects. In the symmetric setting, primeness means that \(M\otimes N\) belonging to the ideal forces \(M\) or \(N\) to belong to it. Cocommutativity of \(A\) supplies the standard symmetry. For a general Hopf algebra, the source invokes noncommutative extensions of this theory without specifying which spectrum or notion of prime tensor ideal is intended.

Write
\[
X_A=\operatorname{Proj}R_A
\]
for the projective cohomological spectrum intended in the source. The cohomology algebra is graded-commutative; the convention for its projective scheme, including treatment of odd degrees and any passage to a commutative even-degree algebra with regrading, must be specified before a scheme-theoretic assertion is made.

The questions attributed to Andruskiewitsch's talk are the following classification and computation problems for such pairs \((k,A)\):
\begin{enumerate}
\item Determine when \(R_A\) is Noetherian. The original excerpt also asks when \(X_A\) is a scheme of finite type, naturally understood over \(k\). Preserve this as an additional question: the supplied parenthetical wording does not establish an equivalence between finite type of \(\operatorname{Proj}R_A\) and Noetherianity of \(R_A\).
\item In a specified tensor triangular framework in which the comparison map is defined and has the indicated target, determine when the source's natural continuous map
\[
\Psi:\operatorname{Spec}(\operatorname{stab}A)\longrightarrow X_A
\]
is a homeomorphism. Its construction, the chosen spectrum for noncocommutative \(A\), and the hypotheses ensuring that its image lies in \(\operatorname{Proj}R_A\) are not supplied in the excerpt and remain qualifications of this question.
\item Determine how to compute \(X_A\) from \(A\), specifying whether the requested output is its underlying topological space or its scheme structure.
\end{enumerate}
The source motivates these questions through classification of thick tensor ideals by specialization-closed subsets and a support criterion for constructing a module from another using shifts, cones, sums, summands, and tensoring with simple modules. Those assertions require the hypotheses of the selected support theory; they are not imposed here as established facts for every finite-dimensional Hopf algebra. The excerpt also mentions \(D^b(A)\) alongside \(\operatorname{stab}A\), but supplies no comparison identifying their spectra. The displayed questions retain the stable-category target of the supplied reconstructed statement.
\subsection{Short English statement}
For a finite-dimensional Hopf algebra over a field, with no cocommutativity assumption, ask when its cohomology is Noetherian, when the stated comparison from the stable-category spectrum to projective cohomology is a homeomorphism, and how to compute that projective spectrum. The original finite-type question is retained separately. Field assumptions, projective-scheme conventions, and the noncommutative spectrum and comparison map require clarification. This is a statement-only draft requiring human mathematical review.
\subsection{Sources}
\begin{itemize}
\item Supplied local UnsolvedMath registry record, permanent ID \(30005780\), title ``Cohomology and Balmer Spectra of Hopf Algebras,'' exact source code \texttt{catalogue entry}. The original excerpt attributes the questions to Andruskiewitsch's talk; the registry supplies no separate proposer.
\item \emph{Mini-Workshop: Bridging Number Theory and Nichols Algebras via Deformations} (2024), Oberwolfach Reports, pp.~235--272, DOI 10.4171/owr/2024/5. Supplied source report and literature-source link labelled ``original OWR report'': \url{https://doi.org/10.4171/owr/2024/5}.
\item Supplied literature-source label ``Nakano--Vashaw--Yakimov, spectrum and support theory'': \url{https://arxiv.org/abs/2112.11170}. No further bibliographic details are supplied.
\item The original excerpt refers to Balmer's results by labels [1] and [2], and to noncommutative extensions by labels [5], [10], and [9]. The corresponding bibliography entries are absent from the supplied record.
\item Supplied literature assessment, checked on \texttt{2026-08-22}: the record labels the problem open and states that finite generation, comparison with the Balmer spectrum, and calculation of projective cohomology are solved only in important special classes, while the three general questions remain open. This is supplied dated curation, not independently verified literature assessment. The background explicitly requests independent verification of time-sensitive status and cited claims.
\item The registry labels the reconstructed statement \texttt{corrected\_verified} and reports that it was checked against the official Oberwolfach report. This is a supplied verification claim; no source visit or independent verification is asserted here.
\end{itemize}
\end{document}
